% Ad Atticum 8.8 (ad-atticum-08-08) --- English translation
% Translated by Claude (Anthropic) from the Latin (Perseus).
% Style guide: STYLE.md.

\ciceroLetterOpener{CICERO ATTICO salutem}

\ciceroSection{1} What a shameful business, and miserable on that very account! For so I feel: that what is miserable, in the end, or rather what alone is miserable, is what is shameful. He had fed Caesar up; the same man he had suddenly begun to fear; he had approved no terms of peace; he had made no preparation for war; he had abandoned the City; he had lost Picenum by his own fault; he had bottled himself up in Apulia; he was setting off for Greece, leaving all of us without a word of farewell \textit{[Greek: aprosphōnētous]}, with no share in a design of his so vast, so unprecedented.

\ciceroSection{2} And now suddenly --- letters from Domitius to him, his own to the consuls. It seemed to me that the honourable \textit{[Greek: to kalon]} had flashed before his eyes, and that he, the man he ought to have been, had cried out: \textit{In the face of this, let them devise what they must and contrive every plot against me; for the right is with me} \textit{[Greek: pros tauth' ho ti chrē kai palamasthōn kai pant' ep' emoi tektainesthōn: to gar eu met' emou]}. But he, bidding the honourable a long farewell \textit{[Greek: polla chairein tō kalō]}, pushes on for Brundisium. Domitius, they say, on hearing the news, has given himself up, and the men who were with him as well. What a thing for mourning! And so grief blocks me from writing you more. I am waiting for your letters.
